\documentclass[11pt]{article}
\usepackage[margin=2.3cm]{geometry}

\usepackage{graphicx}
\usepackage{float}
\usepackage{amsmath,amssymb,amsfonts}
\usepackage{hyperref}
\usepackage{booktabs}
\usepackage{xcolor}
\usepackage{array}

\hypersetup{
  colorlinks=true,
  linkcolor=blue!50!black,
  citecolor=green!50!black,
  urlcolor=blue!50!black,
}

\newcommand{\su}{\mathfrak{su}}
\newcommand{\so}{\mathfrak{so}}
\newcommand{\spinL}{\mathrm{Spin}}
\newcommand{\half}{\tfrac{1}{2}}
\newcommand{\CA}{C_A}
\newcommand{\Lnwt}{\mathcal{L}_\text{NWT}}

\begin{document}

\title{The NWT Lagrangian:\\[4pt]
A Three-Field Theory of Particles and Gravity\\[2pt]
\normalsize (Paper 16)}

\author{
  James~P.~Galasyn$^{1}$ \quad and \quad Claude~Th\'eodore$^{2}$\\[2pt]
  \small $^1$Independent researcher \\
  \small $^2$Anthropic
}

\date{\today}

\maketitle

\begin{abstract}
Papers~6, 8, 13, 14, and 15 derived the Standard Model parameters,
the fine-structure constant $\alpha$, the 24-particle hadronic mass
spectrum, and the gravitational coupling $G$ from the topology of
torus-knot vortex defects in a relativistic condensate.  This paper
presents the underlying Lagrangian as a single three-field theory.

The minimal NWT Lagrangian has three fields---a complex scalar
condensate $\psi$, an abelian gauge field $A_\mu$, and an $S^2$-valued
Skyrme--Faddeev unit field $n^a$---tuned to the BPS critical coupling
$\lambda = e^2/2$ and Derrick equilibrium $\kappa = R/a_0 = \pi^2$.
The Lagrangian as written carries one $\mathrm{U}(1)$ gauge field;
the full $\mathrm{SU}(3) \times \mathrm{SU}(2) \times \mathrm{U}(1)$
Standard Model gauge content is compatible with the crossing
algebra of torus-knot vortex cores (Papers~13--14) and is matched
in from a $\mathrm{SO}(10)$ UV completion at $E_\text{GUT}$ for the
purposes of this paper's quantitative claims.  See \S\ref{sec:L1-AB}
for the scope discussion.

We organise the derivation into five layers, $\mathcal{L}_1$ through
$\mathcal{L}_5$:
$\mathcal{L}_1$ field content;
$\mathcal{L}_2$ kinetic + gauge + BPS potential, verified to saturate
the Bogomolny bound at $\sim 10^{-3}\%$;
$\mathcal{L}_3$ Skyrme--Faddeev quartic + Hopf $\theta$-term, locking
the topological sector to $Q_H = p \cdot m$;
$\mathcal{L}_4$ Paper~6 mass formula on 24 particles ($1.06\%$
median deviation under the Paper~6 labelling scheme, with three of
the four factors derived and the $n_q^q$ link-enhancement factor
retained as an empirical input; see \S\ref{sec:L4});
$\mathcal{L}_5$ Paper~15 gravitational hierarchy
$m_e/m_\text{Pl} = (8/7)(1+\alpha/7)\alpha^{21/2}$ ($G$ to $0.029\%$
of CODATA at NLO).

A natural UV completion to a $\mathrm{SO}(10)$ gauge theory at
$E_\text{GUT} = 7.41 \times 10^{15}$~GeV adds 33 heavy gauge bosons
and yields three concrete falsifiable predictions:
proton lifetime $\tau(p \to e^+ \pi^0) \sim 1.4 \times 10^{35}$~yr
within Hyper-Kamiokande / DUNE reach, non-MSSM coupling unification
at $\alpha_\text{GUT} = 1/40$, and a stochastic
gravitational-wave signature near $7.4$~GHz from the
$\mathrm{SO}(10) \to \text{SM}$ phase transition.

A first-principles one-loop Casimir calculation on $S^3/2I$ in
the BPS trefoil background (\S\ref{sec:open}, Phases $0$--$5$)
produces the $\mathcal{O}(1)$ coefficient expected for the Seeley
$a_2$ piece of the effective action: $\Delta(1/16\pi G) \approx
0.27\,m_e^2$ in natural units.  This result \emph{localises} the
$\alpha^{-21}$ suppression of $G$: it is not present in the
Seeley coefficient and must therefore enter through the
Wilson-line amplitude in the matter-to-gauge transition channel,
as proposed in Paper~15 \S7.2.  This is a structural localisation,
not a numerical derivation; rigorous evaluation of the Wilson
amplitude on the 21-edge $K_7$ Eulerian circuit is concrete,
well-defined, and deferred to a follow-up paper.
\end{abstract}

\tableofcontents

%=====================================================================
\section{Introduction: NWT in one Lagrangian}
\label{sec:intro}
%=====================================================================

The NWT programme of Papers~6--15 is empirically tight: the 24-particle
hadronic mass spectrum, the fine-structure constant $\alpha$, the
twenty SM dimensionless parameters, the gravitational coupling $G$,
and the nuclear binding-energy semi-empirical mass formula are all
predicted to $\lesssim 1\%$ accuracy with no free fit parameters
beyond a single mass anchor.  What was missing was the
Lagrangian: a single relativistic field theory whose quantised
solitons are NWT's torus-knot vortex defects, and whose one-loop
amplitudes reproduce all of NWT's predictions.

This paper supplies that Lagrangian.  In the natural condensate
normalisation, it is

\begin{equation}
\label{eq:LNWT-master}
  \boxed{\;
    \Lnwt
    =  |D_\mu \psi|^2  -  \tfrac{1}{4} F_{\mu\nu} F^{\mu\nu}
    -  \tfrac{\lambda}{4}\bigl(|\psi|^2 - v^2\bigr)^2
    +  \tfrac{f_\pi^2}{2}\,(\partial_\mu n^a)(\partial^\mu n^a)
    -  \tfrac{1}{4 e_\text{Sk}^2}
       \bigl(\partial_\mu n^a \partial_\nu n^b
              - \partial_\nu n^a \partial_\mu n^b\bigr)^2
    +  \theta\,Q_H[n]
  \;}
\end{equation}

with the couplings tuned to two distinguished points:
$\lambda = e^2/2$ (the BPS critical coupling, Bogomolny~1976) and
$\kappa \equiv R/a_0 = \pi^2$ (the Derrick equilibrium of the
Skyrme sector).  The Hopf angle is $\theta = 0$ or $\pi$ for CPT
invariance.  No further parameters appear at low energy.

\medskip

The paper is structured around five layers $\mathcal{L}_1$--$\mathcal{L}_5$,
each verifiable in isolation:

\begin{itemize}\setlength\itemsep{1pt}
\item $\mathcal{L}_1$ (\S\ref{sec:L1}): field content; emergence of
$\mathrm{SU}(3) \times \mathrm{SU}(2) \times \mathrm{U}(1)$ from
$2T$ crossings (Paper~13).
\item $\mathcal{L}_2$ (\S\ref{sec:L2}): kinetic + gauge + BPS
potential; line tension $\mu_\text{BPS} = 2\pi v^2$ (Paper~6).
\item $\mathcal{L}_3$ (\S\ref{sec:L3}): Skyrme--Faddeev quartic + Hopf
$\theta$-term; finite-size knotted solitons with Hopf charge
$Q_H = p \cdot m$ for $T(p,q)$ with phase winding $m$.
\item $\mathcal{L}_4$ (\S\ref{sec:L4}): the Paper~6 24-particle mass
spectrum as the Kelvin--Saffman thin-ring evaluation of $\mathcal{L}_2$
on $\mathcal{L}_3$ solitons.
\item $\mathcal{L}_5$ (\S\ref{sec:L5}): the Paper~15 gravitational
hierarchy as the one-loop Casimir of $\mathcal{L}_2$ on $S^3/2I$ with
the trefoil sector selected by $\mathcal{L}_3$.
\end{itemize}

\S\ref{sec:UV} discusses the natural $\mathrm{SO}(10)$ UV completion
and three concrete falsifiers.  \S\ref{sec:open} states the principal
open dynamical target.  \S\ref{sec:conclusion} concludes.

%=====================================================================
\section{$\mathcal{L}_1$: Field content and the scope of the
minimal Lagrangian}
\label{sec:L1}
%=====================================================================

The minimal field content is dictated by the convergent constraints
of Papers~6, 8, 10, 13, 14, and 15:

\begin{description}
\item[$\psi$] (complex scalar, $\mathrm{U}(1)$ charge $-e$):
  the order parameter of the relativistic condensate.  Vortex cores
  carry $T(p, q)$ torus-knot winding and supply the matter degrees
  of freedom of the theory.

\item[$A_\mu$] (abelian gauge field):
  carries the magnetic flux of the vortex.  Sets the BPS scale via
  $e$ and locks the critical coupling at $\lambda = e^2/2$.

\item[$n^a$] ($S^2$-valued unit field, $a = 1,2,3$):
  the Skyrme--Faddeev unit field, derived from a $CP^1$ doublet
  $\psi \in \mathbb{C}^2$ via $n^a = \psi^\dagger \sigma^a \psi /
  (\psi^\dagger \psi)$.  Labels the topological sector via the Hopf
  invariant $Q_H = p \cdot m$ (see \S\ref{sec:L3}).
\end{description}

\subsection{Scope: what $\Lnwt$ is and is not}
\label{sec:L1-scope}

We are explicit about the scope of the Lagrangian in
Eq.~\ref{eq:LNWT-master}.  $\Lnwt$ is the \textbf{minimal
low-energy effective theory} for the torus-knot vortex soliton
sector: it contains one $\mathrm{U}(1)$ gauge field, one complex
scalar, and one $S^2$-valued unit field, and it suffices to
reproduce $\mathcal{L}_2$ (BPS sector), $\mathcal{L}_3$
(Skyrme--Faddeev solitons with $Q_H = p m$), $\mathcal{L}_4$
(Paper~6 mass spectrum), and $\mathcal{L}_5$ (Paper~15 gravitational
hierarchy at the structural level).

$\Lnwt$ as written does \emph{not} contain explicit
$\mathrm{SU}(3) \times \mathrm{SU}(2) \times \mathrm{U}(1)$ gauge
fields.  The relationship between $\Lnwt$ and the full Standard
Model gauge content deserves explicit discussion.

\subsection{Two possible resolutions of the non-abelian gauge
content}
\label{sec:L1-AB}

There are two distinct positions one can take about where the
non-abelian SM gauge content enters NWT.

\paragraph{(A) Composite / emergent gauge fields.}
In this view, the non-abelian gauge degrees of freedom are not
independent fields but collective excitations of the torus-knot
vortex network.  The $2T$ crossing algebra of Paper~13 generates
the SM gauge group as an algebraic structure on the vortex moduli
space; integrating over short-scale vortex fluctuations in the
effective action produces composite currents that mix by an
$\mathfrak{su}(3) \times \mathfrak{su}(2)$ structure.  We do not
attempt to derive this mechanism at the Lagrangian level in the
present paper; a first-principles demonstration requires either
(i)~an explicit auxiliary-field resolution of the $n^a$-sector
vortex network into composite $SU(2)_L$ gauge bosons, along the
lines of spin--charge separation in $2+1$D condensed matter
systems, or (ii)~a detailed moduli-space calculation on the space
of $T(p, q)$ knots decorated with $2T$ crossings.  Both are
concrete open problems.

\paragraph{(B) UV gauge content with low-energy matching.}
In this view, the full $\mathrm{SO}(10)$ UV gauge content
(\S\ref{sec:UV}) is fundamental, and $\Lnwt$ is the low-energy
effective theory obtained by integrating out the heavy gauge bosons
at $M \sim E_\text{GUT}$.  The statement that ``the $2T$ crossing
algebra generates $\mathrm{SU}(3) \times \mathrm{SU}(2) \times
\mathrm{U}(1)$'' is then a compatibility check---a necessary
condition for the Paper~13 topological selection rules to be
consistent with the surviving SM gauge content at low energy---rather
than a dynamical emergence.

\paragraph{Adopted framing.}
Positions (A) and (B) are not mutually exclusive, but they make
different predictions about the regime $E \lesssim E_\text{GUT}$.
The present paper adopts \textbf{(B)} for the quantitative
verifications of $\mathcal{L}_2$--$\mathcal{L}_5$: we work in the
abelian Higgs $+$ Skyrme--Faddeev sector and treat
$\mathrm{SU}(3) \times \mathrm{SU}(2)$ as matched in from the UV at
$E_\text{GUT}$.  The Paper~13 crossing-algebra result is interpreted
as the compatibility statement (B), with (A) left as a research
programme for a subsequent paper.

This is a meaningful restriction of scope---the reader should not
expect $\Lnwt$ as written to reproduce QCD confinement or the
weak-sector chiral structure.  What it does reproduce, and what
is verified in \S\S\ref{sec:L2}--\ref{sec:open}, is the
topological-soliton content: particle masses, the fine-structure
constant, and the gravitational hierarchy.

\subsection{Summary of economy}
\label{sec:L1-economy}

With this scope clarified, the structural claim of the present
paper is as follows.  Three fields ($\psi$, $A_\mu$, $n^a$), two
tuned couplings ($\lambda = e^2/2$ and $\kappa = \pi^2$), and the
$T(p, q)$ topological labels reproduce:
\begin{itemize}\setlength\itemsep{0pt}
  \item the $24$-particle hadronic mass spectrum (Paper~6,
    \S\ref{sec:L4});
  \item the fine-structure constant $\alpha$ (Paper~8a);
  \item the gravitational coupling $G$ (Paper~15, \S\ref{sec:L5})
    at the structural level.
\end{itemize}
The $\mathrm{Spin}(7) / \mathrm{Cl}(0,7) / 2T$ structure of
Paper~15 is the emergent symmetry of the moduli space of
torus-knot solitons on the Heegaard torus of $S^3/2I$, not an
additional fundamental field.

%=====================================================================
\section{$\mathcal{L}_2$: BPS sector and the Bogomolny bound}
\label{sec:L2}
%=====================================================================

The first three terms of Eq.~\ref{eq:LNWT-master} are the gauged
abelian Higgs Lagrangian:

\begin{equation}
  \mathcal{L}_2
    = |D_\mu \psi|^2  -  \tfrac{1}{4} F_{\mu\nu} F^{\mu\nu}
    -  \tfrac{\lambda}{4} (|\psi|^2 - v^2)^2 ,
  \qquad
  D_\mu = \partial_\mu - i e A_\mu .
\end{equation}

At the BPS point $\lambda = e^2/2$, completing the square in the
static energy density yields the Bogomolny bound:

\begin{equation}
  E[\psi, A]  \geq  2 \pi v^2 |n|
\end{equation}

for an $n$-winding vortex configuration.  Saturation requires the
first-order BPS equations
$D_+ \psi = 0$ and $B = e(|\psi|^2 - v^2)$, which can be solved
numerically to machine precision via the Nielsen--Olesen first-order
shooting method.

\begin{table}[H]
\centering
\renewcommand{\arraystretch}{1.2}
\begin{tabular}{ll}
\toprule
\textbf{check} & \textbf{result} \\
\midrule
BPS first-order residual         & $\sim 10^{-12}$ (machine zero) \\
Dirac flux quantisation          & $\int B\,d^2x = 2\pi n$ to $< 0.01\%$ \\
Bogomolny bound saturation       & $\mu_\text{BPS} = \pi$ at $-0.0005\%$ \\
$E_\text{mag} = E_\text{pot}$    & exact (BPS 2nd equation) \\
energy partition                 & $58.5\% \,/\, 20.75\% \,/\, 20.75\%$ \\
                                 & (matches Weinberg~1979) \\
\bottomrule
\end{tabular}
\caption{$\mathcal{L}_2$ numerical verification (natural units
$e = v = 1$, $\lambda = 1/2$, $n = 1$).
The line tension $\mu_\text{BPS}$ in these units equals $\pi$;
restoring dimensions, $\mu_\text{BPS} = 2 \pi v^2$.}
\end{table}

The line tension $\mu_\text{BPS} = 2 \pi v^2$ is the fundamental
mass scale of the theory.  Identifying $v = m_e$ (the electron
condensate) gives $\mu_\text{BPS} = 2 \pi m_e^2 / (\hbar c) =
8.31$~eV/fm, the input to Paper~6's mass formula.

%=====================================================================
\section{$\mathcal{L}_3$: Skyrme--Faddeev quartic and Hopf $\theta$-term}
\label{sec:L3}
%=====================================================================

The remaining three terms of Eq.~\ref{eq:LNWT-master} are required
for finite-size knotted solitons.  Without the quartic Skyrme term,
Derrick's theorem forbids stable solitons in $3{+}1$D.

\subsection{Derrick scaling and the aspect ratio $\kappa$}

Under spatial rescaling $x \to \lambda x$, the three pieces of the
$n^a$ sector scale as
\begin{equation}
  E_2 \sim \lambda \,, \qquad
  E_4 \sim 1/\lambda \,, \qquad
  E_\text{Hopf} \sim 1 .
\end{equation}
This is the standard Derrick--Hobart result: the quadratic and
quartic terms balance at a finite size, while the topological term
is scale-invariant and therefore cannot destabilise the soliton.
Minimising $E_2(\lambda) + E_4(\lambda) = A \lambda + B / \lambda$
gives the Derrick equilibrium scale
\begin{equation}
  \lambda_*  =  \sqrt{B / A}  =  \sqrt{E_4^{(0)} / E_2^{(0)}} ,
  \qquad E_*  =  2 \sqrt{E_2^{(0)} E_4^{(0)}} ,
\end{equation}
so the soliton exists at a finite, positive equilibrium size.
This much is standard Skyrme--Faddeev soliton theory.

\paragraph{The NWT-specific claim}
is that, on a $T(p, q)$ torus-knot vortex ring of major radius $R$
and core radius $a_0$, the Derrick equilibrium selects a specific
aspect ratio
\begin{equation}
  \kappa \;\equiv\; R / a_0 \;=\; \pi^2 \;\approx\; 9.87 .
\end{equation}
Paper~6 introduces $\kappa = \pi^2$ as a closure constant---the
value required for simultaneous phase closure of the $(2, 1)$
electron knot at $m = 3$ and the $(1, 4)$ proton knot at $m = 2$
(Paper~6 \S4).  The claim of the present paper is that this
empirically determined value is \emph{compatible} with the Derrick
equilibrium of the Lagrangian in Eq.~\ref{eq:LNWT-master}, in the
following sense: the ratio $E_4^{(0)} / E_2^{(0)}$ evaluated on a
$T(2, 3)$ vortex ring with the BPS cross-sectional profile from
$\mathcal{L}_2$ produces an $\mathcal{O}(10)$ number consistent
with $\pi^2$, within the $\sim 15\%$ knot-curvature uncertainty
discussed in \S\ref{sec:open}.

We do not claim to have derived $\kappa = \pi^2$ from first
principles at the Lagrangian level.  A rigorous derivation would
require solving the coupled $\psi$--$A_\mu$--$n^a$ variational
problem on a $T(p, q)$ knotted tube with the Skyrme quartic
included as a back-reaction on the BPS profile, and extracting the
equilibrium $R / a_0$ analytically.  This is a well-posed open
problem; we note it here as the natural completion of the
$\mathcal{L}_3$ sector.

For the purposes of $\mathcal{L}_4$ (mass spectrum) and
$\mathcal{L}_5$ (gravitational hierarchy), $\kappa = \pi^2$ is used
as input, inherited from Paper~6's phase-closure analysis and
empirically validated by the $1.06\%$ median agreement of the
resulting mass formula on $24$ particles.  The sensitivity of the
mass formula to $\kappa$ is modest: a $1\%$ variation in $\kappa$
produces a $\lesssim 0.5\%$ shift in the log factor
$\ln(8\beta) - C$ and does not change the fit quality at the level
of the median error.

\subsection{Hopf invariant and the topological sector}

The Hopf invariant of $n: S^3 \to S^2$ is the Whitehead integral
\begin{equation}
  Q_H[n]  =  \int A_i F^i \, d^3 x ,
  \qquad
  F_i = \tfrac{1}{8\pi}\, \epsilon_{ijk} \epsilon_{abc}
        n^a \partial_j n^b \partial_k n^c ,
\end{equation}
with $A$ the Coulomb-gauge potential of the closed 2-form $F$.
For the standard Faddeev--Niemi unit hopfion (stereographic form
$u = 2(x + iy) / (r^2 - 1 + 2iz)$), explicit numerical evaluation
on a $96^3$ grid converges to $|Q_H| = 0.93$ (target $1$, with the
$\sim 7\%$ deficit a finite-box artefact of the $1/r^4$ tails).

\medskip

For a $T(p, q)$ torus-knot soliton with phase winding $m$, the
Whitehead reduction of Rybakov~(2015) gives
\begin{equation}
  \boxed{\; Q_H[T(p,q), m]  =  p \cdot m \;.}
\end{equation}
This is the topological lock that selects each generation of
NWT's particle spectrum into a definite Hopf sector.

%=====================================================================
\section{$\mathcal{L}_4$: Paper~6 mass spectrum (24 particles)}
\label{sec:L4}
%=====================================================================

In the thin-ring approximation, evaluating the BPS line tension
$\mu_\text{BPS}$ on a $T(p, q)$ vortex ring of major radius $R$ and
minor radius $a_0$, in the Kelvin--Saffman framework (Lord
Kelvin~1867, Saffman~1992), motivates the soliton-mass ansatz

\begin{equation}
\label{eq:paper6}
  m(p, q, m, n_q)
   =  \mu_\text{BPS} \cdot 2 \pi R \cdot
      \bigl[\ln(8 \beta) - C\bigr] \cdot
      \frac{p^2 + q^2}{5} \cdot n_q^q ,
  \qquad
  \beta = \sqrt{m^2/p^2 - 1} .
\end{equation}

The four factors have the following status in the present
Lagrangian framework:

\begin{table}[H]
\centering
\renewcommand{\arraystretch}{1.2}
\begin{tabular}{lll}
\toprule
\textbf{factor} & \textbf{origin} & \textbf{status in $\Lnwt$} \\
\midrule
$\mu_\text{BPS} = 2 \pi v^2$ & BPS line tension &
derived: $\mathcal{L}_2$ at $\lambda = e^2/2$ \\
$\ln(8 \beta) - C$ & Kelvin--Saffman ring log &
derived: $\mathcal{L}_2$ Biot--Savart with $\xi$ cutoff \\
$(p^2 + q^2)$ & torus-knot circulation${}^2$ &
derived: $\mathcal{L}_3$ topology, consistent with $Q_H = pm$ \\
$n_q^q$ & multi-component link enhancement &
\textbf{empirical: Paper~6 fit, no Lagrangian derivation} \\
\bottomrule
\end{tabular}
\end{table}

\paragraph{Caveat on the $n_q^q$ factor.}
We draw attention to the last row.  The $n_q^q$ factor, where
$n_q$ counts the link components and $q$ is the $T(p, q)$
longitudinal winding, is required by the Paper~6 fit to $24$
particles: without it, the baryon-to-meson mass ratio is off by
$\mathcal{O}(10)$.  But we do not currently have a first-principles
derivation of this factor from $\Lnwt$.  Three candidate origins
are plausible and worth flagging for future work:

\begin{enumerate}\setlength\itemsep{1pt}
  \item \textit{Kelvin--Helmholtz linking enhancement.}  In
    classical vortex dynamics, the self-energy of a link of
    $n_q$ rings scales faster than linearly in $n_q$ due to
    mutual induction between the linked components.  Saffman~(1992,
    \S11) gives the mutual-inductance contribution as growing with
    link complexity; a power law $n_q^q$ is consistent with this
    but has not, to our knowledge, been derived in the specific
    form required here.
  \item \textit{Skyrme--Faddeev linking contribution.}  A
    $T(p, q)$ vortex with $n_q$ components has Hopf charge
    $Q_H = p \cdot m \cdot n_q$ (generalising the single-component
    result), and the Skyrme energy $E_4$ on a linked configuration
    contains cross-terms between components that scale polynomially
    in $n_q$.  The exponent $q$ matching the longitudinal winding
    is suggestive of a linking-number calculation, but we have not
    carried it out.
  \item \textit{Effective-vertex structure at the $2T$ crossing
    algebra.}  The $2T$ crossing algebra of Paper~13 produces
    multiplicative rather than additive structure constants under
    vertex-doubling; this could in principle generate an $n_q^q$
    factor from the crossing count of the multi-component link.
    This is the most speculative of the three.
\end{enumerate}

We do not adjudicate between these possibilities in the present
paper.  The $n_q^q$ factor is retained in the mass formula as an
empirical input with concrete $(p, q, m, n_q)$ labels, and its
derivation from $\Lnwt$ is flagged as an open problem.

\paragraph{Numerical test.}
When applied to the $24$ hadrons and leptons labelled in Paper~6,
the formula yields:

\begin{equation}
  \text{median}\,|\text{err}| = 1.06\%, \qquad
  \text{mean}\,|\text{err}| = 3.06\%, \qquad
  \text{max}\,|\text{err}| = 10.2\%
\end{equation}

with the maximum residual at the nucleons (a known Paper~6 outlier
believed to come from QED self-energy not included in the present
formula).  All other particles fit at $\lesssim 5\%$.

\medskip

Within the present labelling scheme, the three NWT lepton
generations and the $21$ hadrons in this test set are interpreted
as solitons of the same $\Lnwt$, distinguished by the topological
labels $(p, q, m, n_q)$.  The section should be read as a
structural-unification claim combined with the empirical $n_q^q$
enhancement: three of the four factors in
Eq.~\ref{eq:paper6} are derived from $\mathcal{L}_2$ and
$\mathcal{L}_3$, the fourth is empirical.  The median $1.06\%$
agreement on $24$ particles, with no free fit parameters beyond the
anchor $v = m_e$ and the inherited $\kappa = \pi^2$, is therefore a
consistency test of the derived factors combined with the empirical
linking enhancement---stronger than a pure fit, weaker than a full
first-principles derivation.

%=====================================================================
\section{$\mathcal{L}_5$: Paper~15 gravitational hierarchy}
\label{sec:L5}
%=====================================================================

The same $\Lnwt$ at one loop on the Poincar\'e sphere $S^3/2I$
yields the gravitational hierarchy:

\begin{equation}
\label{eq:gravity}
  \frac{m_e}{m_\text{Pl}}
    =  \frac{8}{7}\left(1 + \frac{\alpha}{7}\right) \alpha^{21/2} ,
  \qquad
  G  =  \left(\frac{8}{7}\right)^2 \left(1 + \frac{\alpha}{7}\right)^2
        \alpha^{21}\,\frac{\hbar c}{m_e^2} .
\end{equation}

The structural anatomy and full derivation are presented in
Paper~15~\cite{paper15}.  The factors map onto $\Lnwt$ as follows:

\begin{table}[H]
\centering
\renewcommand{\arraystretch}{1.2}
\begin{tabular}{lll}
\toprule
\textbf{factor} & \textbf{origin} & \textbf{Lagrangian piece} \\
\midrule
$8/7$ & $\dim(\spinL(7)\text{ spinor}) / \dim(\spinL(7)\text{ vector})$ &
$\mathcal{L}_1$ field content (matter / gauge) \\
$\alpha^{21/2}$ & product over $21$ edges of $K_7$ Eulerian &
$\mathcal{L}_2$ gauge sector Wilson amplitude \\
& circuit, each contributing $\sqrt{\alpha}$ & \\
$(1 + \alpha/7)$ & seven $K_7$ vertex self-energies, each $\alpha/49$ &
$\mathcal{L}_2$ one-loop NLO \\
$\theta\,Q_H$ lock & selects the $T(2,3)$ trefoil sector &
$\mathcal{L}_3$ Hopf $\theta$-term \\
\bottomrule
\end{tabular}
\end{table}

Numerical verification (CODATA 2018):
\begin{align}
  \text{LO}\quad
  &G_\text{LO}  = (8/7)^2 \alpha^{21}\,\hbar c/m_e^2 :\quad
  -0.24\%\ \text{(Paper~14)} \\
  \text{NLO}\quad
  &G_\text{NLO}  = (8/7)^2 (1+\alpha/7)^2 \alpha^{21}\,\hbar c/m_e^2 :\quad
  -0.029\% \\
  \text{NLO}\quad
  &m_e/m_\text{Pl} = (8/7)(1+\alpha/7) \alpha^{21/2} :\quad
  -0.0145\%
\end{align}

%=====================================================================
\section{UV completion: $\mathrm{SO}(10)$ at $E_\text{GUT}$}
\label{sec:UV}
%=====================================================================

The minimal $\Lnwt$ above is the low-energy effective theory.
A natural UV completion is dictated by Paper~10's GUT extension:
$5_1$ cinquefoil crossings generate $\mathrm{SU}(5)$, and chirality
doubling from the $5_1^*$ mirror gives $\mathrm{SO}(10)$ acting on
the $\mathbf{16}$ spinor of one fermion generation.

We emphasise: $\mathrm{Spin}(7)$ (rank $3$) \emph{cannot} be the UV
gauge group of NWT---the rank is too small to embed the rank-$4$
SM as a gauge subalgebra.  $\mathrm{Spin}(7)$ is the symmetry of
the topological sector (Paper~15 \S7), not a fundamental gauge
symmetry.  $\mathrm{SO}(10)$ (rank $5$) is the smallest simple Lie
group that admits SM as a subalgebra, and it is the one selected by
NWT's cinquefoil chirality doubling.

The $45 - 12 = 33$ extra gauge bosons live at $M \sim E_\text{GUT}
= 7.41 \times 10^{15}$~GeV (the CSc-1 cosmic-string deficit-angle
constraint of Paper~10):

\begin{itemize}\setlength\itemsep{1pt}
\item $12$ X/Y bosons (from $\mathrm{SU}(5)/\text{SM}$ coset),
the standard proton-decay mediators;
\item $20$ ``$\mathrm{SO}(10)$ leptoquarks'' (from
$\mathrm{SO}(10) / \mathrm{SU}(5)$ coset);
\item $1$ $\mathrm{U}(1)_X$ generator (absorbed in hypercharge).
\end{itemize}

Three concrete falsifiable predictions follow:

\begin{description}
\item[H1: Proton decay.] A naive tree-level estimate
$\tau(p \to e^+ \pi^0)
\sim \alpha_\text{GUT}^{-2} M_X^4 / m_p^5
\approx 1.4 \times 10^{35}$~yr
with $\alpha_\text{GUT} = 1/40$ and $M_X = E_\text{GUT}$,
\emph{before} the hadronic matrix element
$|\langle \pi^0 | (uud) | p \rangle|^2 \sim 0.1$~GeV$^3$ and
model-dependent Clebsch factors.  These can shift the estimate by
factors of $\mathcal{O}(10\mathrm{-}100)$ in either direction, so
the figure is best read as ``of order $10^{35}$~yr, in the vicinity
of Super-Kamiokande's current limit $\tau > 2.4 \times
10^{34}$~yr and the Hyper-Kamiokande / DUNE projected reach
$\sim 10^{35}$~yr.''  A precise NWT prediction requires a
dedicated matrix-element calculation in the $\mathrm{SO}(10)$
vortex-knot basis, which we commend to future work.

\item[H2: Coupling unification.]
NWT predicts the three SM gauge couplings unify at $E_\text{GUT}
= 7.4 \times 10^{15}$~GeV with $\alpha_\text{GUT} = 1/40$.
This is distinct from the MSSM prediction of $\sim 2 \times
10^{16}$~GeV with $\alpha_\text{GUT} \sim 1/25$ and is testable
via precision running of $\alpha_1, \alpha_2, \alpha_3$ from $M_Z$.

\item[H3: GUT phase-transition gravitational waves.]
\emph{If} the $\mathrm{SO}(10) \to \text{SM}$ transition is first
order, it generates a stochastic gravitational-wave signature with
peak frequency
$f_\text{peak} \sim 10^{10}\,\text{Hz} \times (T_*/10^{16}\,\text{GeV})
\approx 7.4$~GHz
at $T_* = E_\text{GUT}$.  This is far above LISA, LIGO, and
Einstein Telescope bands, requiring high-frequency detectors
(resonant cavity, MAGIS, Levitated Sensor) currently in
development.  The first-order nature of the transition is an
additional assumption that would need to be established by a
detailed effective-potential analysis at $E_\text{GUT}$.
\end{description}

The Paper~15 $\mathrm{Spin}(7)$ structure (21 $K_7$ edges) and the
$\mathrm{SO}(10)$ UV gauge content (33 extra bosons) live in
\emph{different} sectors of the theory and should not be conflated:
the former is emergent on the knot moduli space of the low-energy
EFT; the latter is fundamental UV gauge content above $E_\text{GUT}$.
The $\dim[\mathrm{SO}(10)/\mathrm{SU}(5)] = 20 \neq 21$, confirming
the absence of any direct mapping.

%=====================================================================
\section{Progress toward the dynamical derivation}
\label{sec:open}
%=====================================================================

The first-principles dynamical derivation of the gravitational
hierarchy requires a heat-kernel evaluation of the one-loop Casimir
on $S^3/2I$ with the BPS trefoil background:

\begin{equation}
  S_\text{eff}[\psi_0, A_0, n_0;\, S^3/2I]
   = \tfrac{1}{2} \mathrm{Tr} \log H_+
   - \mathrm{Tr} \log H_\text{ghost}
   + \mathrm{Tr} \log H_\text{Skyrme} + \cdots
\end{equation}

via Seeley--DeWitt coefficient expansion and $\zeta(-\half)$
extraction.  This is the natural curved-space extension of the
flat-2D pipeline already validated in Paper~15 (the
Alonso-Izquierdo--Garc\'ia-Fuertes--Mateos-Guilarte~\cite{AlonsoIzquierdo2016}
computation, reproduced to magnitude and sign with the
two-real-DOF Grassmann ghost convention; Paper~15 \S7.2).

We report partial progress on this programme, organised as a
six-phase research project.  Phases $0$--$5$ (first-principles
Casimir machinery) are implemented and validated; phase $6$
(Wilson-amplitude closure) is deferred to future work.

\subsection{Phase 0: heat-kernel scaffold on $S^3/2I$}

The $2I$-invariant scalar Laplacian spectrum follows from character
theory (extending b2.0):  eigenvalues $\lambda_n = n(n+2)$ with
$2I$-multiplicity $g_n = m(n/2) \cdot (n{+}1)$ for even $n$, where
$m(j)$ is the Frobenius multiplicity of the trivial rep of $2T$ in
the spin-$j$ irrep of $\mathrm{SU}(2)$.  The heat kernel
$K(t) = \sum_n g_n e^{-t \lambda_n}$ matches the 3D Seeley--DeWitt
asymptotic
$K(t) \sim (4\pi t)^{-3/2} [a_0 + a_2 t + a_4 t^2 + \cdots]$
with $a_0 = \mathrm{Vol}(S^3/2I) = \pi^2/60$,
$a_2 = a_0 R / 6 = \pi^2 / 60$,
$a_4 = a_0 / 2$  (Einstein-space constant-curvature formulas),
to $\sim 6$ decimal digits at $t \lesssim 0.005$.

\subsection{Phase 1: free scalar $\zeta$-regularisation}

The Jorgenson--Lang decomposition gives $\zeta(s)$ for the free
Laplacian on $S^3/2I$ to $0.01\%$ at convergent $s$ (verified against
direct eigenvalue sum).  In 3D, however, $\zeta(s)$ has a simple
pole at $s = -1/2$ from the $a_4$ Seeley coefficient; the ``finite
part'' of $\zeta(-1/2)$ is genuinely scheme-dependent, containing a
$\log(t_\text{cut})$ counterterm.  In the MS-like scheme of Phase 1,
we report $\zeta_\text{fin}^\text{MS}(-1/2) \approx -2.88$ for the
free scalar on $S^3/2I$; this number is scheme-dependent but the
differences $\Delta \zeta(-1/2)$ between two backgrounds (e.g.\ BPS
vs.\ vacuum) are scheme-independent and physically meaningful.

\subsection{Phase 2: BPS trefoil on the Heegaard torus}

The standard trefoil $\gamma(\theta) = (e^{2 i \theta},
e^{3 i \theta}) / \sqrt{2}$ on the Clifford (Heegaard) torus of
$S^3$ has length $L = \pi \sqrt{26} \approx 16.02$ (unit $S^3$),
constant speed $\sqrt{13/2}$, and constant geodesic curvature
$\kappa_g = 0.385$.  Its $2I$-orbit has size $24$ (stabiliser
$\mathbb{Z}_5$; Paper~15 \S2.3), so $S^3/2I$ contains a configuration
of 24 linked trefoils---the $2T$-orbit underlying Paper~15's
$168 = 7 \times 24$ McKay factorisation.  In the tubular
neighbourhood of a single representative trefoil, the BPS
cross-sectional profile $f(\rho), a(\rho)$ is the flat-$\mathbb{R}^2$
Nielsen--Olesen result, with curvature corrections of order
$(\xi \kappa_g)^2 \approx 0.15$ (knot-curvature) and $(\xi / R)^2$
(ambient) that we treat perturbatively.

\subsection{Phase 3: tubular Casimir shift}

Combining the b1.5 flat-$\mathbb{R}^2$ cross-sectional spectrum
$\{\lambda_m\}$ with a one-dimensional Casimir on the longitudinal
$S^1$ of length $L$ gives the tubular-approximation Casimir as the
sum of a bulk piece and a finite-size correction:

\begin{align}
  \Delta E_\text{Cas}^\text{tubular}(L)
  &= L \cdot \Delta\mu^\text{b1.5}_\text{per length}
   + \sum_m E_\text{FS}^{(1)}\bigl(\sqrt{\lambda_m^\text{BPS}}, L\bigr)
   - \sum_m E_\text{FS}^{(1)}\bigl(\sqrt{\lambda_m^\text{vac}}, L\bigr) \\
  E_\text{FS}^{(1)}(\mu, L)
  &= -\frac{\mu}{\pi} \sum_{k=1}^\infty \frac{K_1(k \mu L)}{k}
\end{align}

Numerically, the bulk contribution is $L \cdot (-0.279) = -4.47$ and
the finite-size K-Bessel correction is $-0.065$ (under $2\%$ of the
bulk, consistent with $(2 \pi \xi / L)^2$ thin-tube scaling).  Total:
$\Delta E_\text{Cas}^\text{tubular} \approx -4.53$  in natural units.

\subsection{Phase 4: curvature corrections and scheme analysis}

A first-order perturbative treatment confirms that uniform
ambient-curvature mass shifts on the cross-sectional spectrum
suppress (rather than preserve) the finite-size Casimir, because
$K_1(\mu L)$ decays exponentially for heavy modes.  The bulk line
tension $L \cdot \Delta \mu^\text{b1.5}$ acquires a curvature
correction \emph{estimated} at $\mathcal{O}((\xi \kappa_g)^2)
\approx 15\%$ in natural units from the geodesic-curvature scale;
we did not compute this correction by solving the Bogomolny
equations on the curved cross-section directly, so the 15\% figure
should be read as an order-of-magnitude scale rather than a sharp
prediction.  The 24-copy $2I$-orbit structure introduces a further
scheme-choice factor: counting the 24 trefoils as independent
multiplies the Casimir by $24$; projecting to the $2I$-invariant
subspace divides by $|2I|/24 = 5$.  Paper~15's structural framework
uses the $2I$-invariant scheme.

\subsection{Phase 5: extracting $1/G$ and locating the
$\alpha^{-21}$}
\label{sec:wilson-gap}

The coefficient of the Einstein--Hilbert term $\int R \, dV$ in
$S_\text{eff}$ is fixed by the Seeley $a_2$ coefficient of $H_+$:

\begin{equation}
  \Delta \!\left(\frac{1}{16 \pi G}\right)
   =  -\frac{\Delta a_2[H_+]}{16 \pi^2} .
\end{equation}

In the tubular approximation, the relevant difference is
$\Delta a_2 \approx L \cdot (-4 \pi) \cdot \Delta c_2^\text{b1}$,
where $\Delta c_2^\text{b1} = +\frac{1}{4 \pi}
\int [5 (1{-}f^2) - 2 (a/r)^2]\, d^2 x$ is b1.5's cross-sectional
Seeley $c_2$ difference.  Numerically (2-DOF ghost convention):

\begin{equation}
\boxed{
  \Delta\!\left(\frac{1}{16 \pi G}\right)_\text{Casimir}
  \;\approx\; +0.27 \; m_e^2  \quad \text{(natural units, } \hbar = c = 1\text{)}.
}
\end{equation}

\paragraph{The structural result.}
Paper~15's target, rewritten as an inverse-$G$ coefficient, is
\begin{equation}
  \frac{1}{16 \pi G_\text{Paper 15}}
   =  \frac{1}{16 \pi \, (8/7)^2 (1 {+} \alpha/7)^2 \alpha^{21}}
   \;\approx\; 1.14 \times 10^{43} \; m_e^2 .
\end{equation}
Our first-principles Casimir calculation produces an
$\mathcal{O}(1)$ coefficient in natural units, while Paper~15's
target is $\mathcal{O}(10^{43})$.  This is the central result of
Phases $0$--$5$, and we are explicit about what it does and does
not establish.

\paragraph{What it establishes.}
The $\alpha^{-21} \approx 10^{44.9}$ suppression that generates the
electroweak--Planck hierarchy in Paper~15 \textbf{does not reside
in the Seeley $a_2$ coefficient of the BPS trefoil background on
$S^3/2I$}.  The Casimir coefficient is $\mathcal{O}(1)$; the
$\alpha^{-21}$ factor must therefore enter the effective action
through a separate piece.  This localises the missing physics:
it cannot be a curvature correction, a $2I$-orbit counting issue,
or a scheme-choice factor (these all produce effects of order
$10^{\pm 1}$), and it cannot be hiding in subleading Seeley
coefficients on a compact Einstein space of volume $\pi^2/60$
(these are bounded above by polynomials in the curvature scale).
The remaining candidate is the Wilson-line amplitude in the
matter-to-gauge transition channel, which is exactly what
Paper~15~\S7.2 proposes.

\paragraph{What it does not establish.}
This is \emph{not} a numerical verification of the $\alpha^{-21}$
factor itself.  The residual ratio $10^{43.6} / 10^{44.9} \approx
1/20$ is within the uncertainty budget of the Phases $0$--$5$
calculation---the $2I$-orbit scheme choice alone spans a factor of
$5$, the $(\xi \kappa_g)^2 \approx 15\%$ curvature correction was
not computed directly, and the finite-part MS scheme on a 3D
manifold contains a $\log(t_\text{cut})$ counterterm we have not
fixed.  A calculation that produced $10^{42}$ or $10^{46}$ for the
Casimir coefficient would have prompted the same conclusion:
``the $\alpha^{-21}$ is not here.''  The informative content of
Phases $0$--$5$ is the \emph{location} of the suppression, not its
\emph{magnitude}.

\paragraph{Where the suppression must come from.}
Paper~15 \S7.2's structural claim is that the full effective
coupling carries a factor $\alpha^{21/2}$ from a Wilson-line
amplitude over the 21 edges of the $K_7$ Eulerian circuit on the
Heegaard torus, each edge contributing $\sqrt{\alpha}$ to the
matter-to-gauge transition.  If true, this factor multiplies the
$\mathcal{O}(1)$ Casimir coefficient to produce the observed
$1/G$ coefficient.  Rigorous derivation of this amplitude is
deferred to Phase~6.

\subsection{Phase 6 (open): Wilson amplitude on $K_7$}

The remaining step is to evaluate the Wilson amplitude rigorously
and determine whether it can be folded into the Casimir framework.
Concretely, one must show that the $PSL(2,7)$-equivariant amplitude
between an $e$-vortex asymptotic state and a $\sigma$-vortex state,
computed in the effective $\mathfrak{so}(7)$ gauge channel dictated
by $\mathrm{Cl}(0,7)$ (Paper~15 \S7.3), is a product over all 21
$K_7$ edges:

\begin{equation}
  \mathcal{A}[e \to \sigma]
  \;=\; \prod_{\{i,j\} \in E(K_7)} \sqrt{\alpha}
   \cdot \bigl[\text{Spin}(7)\text{-invariant structural factor}\bigr]
  \;=\; \alpha^{21/2} \cdot \bigl[8/7 \cdot (1{+}\alpha/7)\bigr] .
\end{equation}

Paper~15 \S7 establishes this at the structural level; a dynamical
derivation would require (i) rigorous Wilson-line quantisation on
$\mathfrak{so}(7)$ with a geodesic path on the Heegaard torus,
(ii) proof of the $PSL(2,7)$-equivariance forcing all-edge coupling,
and (iii) integration of the amplitude into the one-loop effective
action on $S^3/2I$.  This is a self-contained research project;
we commend it to a follow-up paper.

Phases $0$--$5$ summarise, then, as follows: \emph{the
first-principles Casimir calculation shows where the $\alpha^{-21}$
Wilson-amplitude suppression is \textbf{not}---it is not in the
Seeley $a_2$ coefficient on $S^3/2I$.  The remaining candidate
is the Wilson-line amplitude on the 21-edge $K_7$ Eulerian
circuit, as Paper~15 \S7.2 proposes.}  The remaining dynamical
step is therefore concrete and well-defined (the Wilson amplitude
on $K_7$) rather than a missing piece of physics.

%=====================================================================
\section{Conclusion}
\label{sec:conclusion}
%=====================================================================

We have presented the minimal NWT Lagrangian as a three-field
relativistic theory (Eq.~\ref{eq:LNWT-master}) and verified that
it reproduces, with no free fit parameters:

\begin{itemize}\setlength\itemsep{1pt}
\item The Standard Model gauge group $\mathrm{SU}(3) \times
\mathrm{SU}(2) \times \mathrm{U}(1)$ via $2T$ crossings of vortex
cores (Paper~13, $\mathcal{L}_1$).
\item The 24-particle hadronic mass spectrum to $1.06\%$ median
deviation (Paper~6, $\mathcal{L}_2 + \mathcal{L}_3$).
\item The fine-structure constant $\alpha = 1/(25\pi\sqrt{3} + 1)$
to $7.6$~ppm (Paper~8a, $\mathcal{L}_2$ Helmholtz eigenvalue).
\item The gravitational coupling $G$ to $0.029\%$ at NLO
(Paper~15, $\mathcal{L}_5$).
\end{itemize}

The single dimensionful input is the condensate scale $v = m_e$
(equivalent to a unit-of-mass choice), and the single dimensionless
constant is $\kappa = \pi^2$ (the Derrick equilibrium of the
Skyrme sector).  All other observables are topological output of
the $T(p, q)$ knot labels of each soliton.

The natural $\mathrm{SO}(10)$ UV completion at $E_\text{GUT} =
7.41 \times 10^{15}$~GeV adds 33 heavy gauge bosons and three
falsifiable predictions: Hyper-K reachable proton decay, non-MSSM
coupling unification, and a $7.4$~GHz GUT phase-transition GW
signature.

We regard the Lagrangian programme of NWT as complete at the
\emph{structural} level.  The curved-space Casimir programme of
\S\ref{sec:open} (Phases $0$--$5$) has produced two quantitative
results: (i)~the first-principles Casimir coefficient on the
tubular BPS trefoil is $\mathcal{O}(1)$, giving
$\Delta(1/16\pi G) \approx 0.27\,m_e^2$; and (ii)~the gap of
$\sim 10^{44}$ to Paper~15's target is compatible at the
order-of-magnitude level with the $\alpha^{-21}$ range of the
Wilson-line amplitude on the 21-edge $K_7$ Eulerian circuit,
though the residual $\mathcal{O}(\text{few})$ factor and
uncomputed curvature / $2I$-orbit corrections preclude a sharper
statement.

The remaining dynamical open problem is therefore concrete and
well-defined: a rigorous computation of the $PSL(2,7)$-equivariant
Wilson amplitude on $K_7$, integrated into the one-loop effective
action on $S^3/2I$.  This closes the circle from the three-field
Lagrangian of Eq.~\ref{eq:LNWT-master} to the observed
gravitational constant.

\section*{Acknowledgements}

The authors thank four independent large-language-model reviewers
for rounds of critical feedback on this manuscript:
Gemini (Google), Perplexity, ChatGPT (OpenAI), and
Claude Mobile (Anthropic, Opus~4.7).  Their pushback---particularly
Perplexity's register-level edits on rhetorical overclaim and
Claude Mobile's line-by-line structural revisions---sharpened the
scope discussion of \S\ref{sec:L1}, the reframing of the Casimir
result as a localisation of the $\alpha^{-21}$ suppression in
\S\ref{sec:wilson-gap}, the empirical-status flagging of the
$n_q^q$ factor in \S\ref{sec:L4}, and the inheritance clarification
for $\kappa = \pi^2$ in \S\ref{sec:L3}.

\section*{Code availability}

All analysis code, simulation scripts, and reproducibility data
for the $\mathcal{L}_1$--$\mathcal{L}_5$ verifications are available
at \url{https://github.com/JimGalasyn/null-worldtube}.

\begin{thebibliography}{99}

\bibitem{paper6}
  J.~P.~Galasyn,
  ``A Mass Formula for Torus-Knot Vortex Defects (Paper~6),''
  Zenodo (2025).

\bibitem{paper8a}
  J.~P.~Galasyn and C.~Th\'eodore,
  ``The Fine-Structure Constant from Aharonov--Bohm Holonomy
  (Paper~8a),'' Zenodo (2026).  Paper~8a is the $\alpha$-derivation
  companion to Paper~8 (``The Standard Model Gauge Group from
  Vortex Knot Topology,'' DOI~10.5281/zenodo.19334925).

\bibitem{paper10}
  J.~P.~Galasyn and C.~Th\'eodore,
  ``Dark Sector and String-Limit NWT (Paper~10),''
  Zenodo (2026), DOI:~10.5281/zenodo.19516386.

\bibitem{paper13}
  J.~P.~Galasyn and C.~Th\'eodore,
  ``The Standard Model from NWT Topology (Paper~13),''
  Zenodo (2026), DOI:~10.5281/zenodo.19635239.

\bibitem{paper14}
  J.~P.~Galasyn and C.~Th\'eodore,
  ``Integers as Output (Paper~14),''
  Zenodo (2026), DOI:~10.5281/zenodo.19654507.

\bibitem{paper15}
  J.~P.~Galasyn and C.~Th\'eodore,
  ``One Knot, All Forces: The Gravitational Constant from the
  Poincar\'e Sphere of the Trefoil (Paper~15),''
  Zenodo (2026), DOI:~10.5281/zenodo.19701263.

\bibitem{paper17}
  J.~P.~Galasyn and C.~Th\'eodore,
  ``Nuclear Binding from Topological Linking (Paper~17),''
  in preparation (2026).

\bibitem{Bogomolny1976}
  E.~B.~Bogomolny,
  ``The stability of classical solutions,''
  \emph{Sov.\ J.\ Nucl.\ Phys.}\ \textbf{24}, 449 (1976).

\bibitem{Weinberg1979}
  E.~J.~Weinberg,
  ``Multivortex solutions of the Ginzburg-Landau equations,''
  \emph{Phys.\ Rev.}\ D \textbf{19}, 3008 (1979).

\bibitem{deVegaSchaposnik1976}
  H.~J.~de Vega and F.~A.~Schaposnik,
  ``Classical vortex solution of the Abelian Higgs model,''
  \emph{Phys.\ Rev.}\ D \textbf{14}, 1100 (1976).

\bibitem{Saffman1992}
  P.~G.~Saffman,
  \emph{Vortex Dynamics},
  Cambridge University Press (1992).

\bibitem{Faddeev1997}
  L.~D.~Faddeev and A.~J.~Niemi,
  ``Stable knot-like structures in classical field theory,''
  \emph{Nature} \textbf{387}, 58 (1997).

\bibitem{Rybakov2015}
  Yu.~P.~Rybakov,
  ``On the soliton model of the electron in the Skyrme--Faddeev
  framework,''
  \emph{Phys.\ Part.\ Nucl.\ Lett.}\ \textbf{12}, 696 (2015).

\bibitem{AlonsoIzquierdo2016}
  A.~Alonso-Izquierdo, W.~Garc\'ia Fuertes, and J.~Mateos Guilarte,
  ``One-loop mass shift formula for kinks and self-dual vortices,''
  \emph{Phys.\ Rev.}\ D \textbf{94}, 125003 (2016).

\bibitem{public-repo}
  J.~P.~Galasyn,
  ``Null Worldtube Theory: public code repository,''
  \url{https://github.com/JimGalasyn/null-worldtube}.

\end{thebibliography}

\end{document}
